% !PRINTED 679
\weberpage{679}
For if we denote by $x_i'$ integral rational functions of $z'$, of degree at most $m$, then
\weq{7}{x_i = z^m x_i'}
is an integral rational function of $z$. Now if we had
\[
\lambda_s' = \sum_{0,\,s-1}^{i} x_i'\lambda_i',
\]
then from (6) and (7) it would follow:
\[
z^m\lambda_s = \sum_{0,\,s-1}^{i} x_i z^{r_s-r_i}\lambda_i,
\]
and this is impossible, because the $\lambda_i$ are linearly independent; and there can also be no function $\lambda_s'$ of lower exponent than $r_s$, because from it one could derive a function $\lambda_s$ of lower exponent.

The field discriminants $D_z, D_{z'}$ with respect to the variable $z$ and $z'$ are
\[
\begin{aligned}
D_z &= \Delta(\lambda_1, \lambda_2, \ldots, \lambda_n), \\
D_{z'} &= \Delta(\lambda_1', \lambda_2', \ldots, \lambda_n'),
\end{aligned}
\]
and from this it follows by (6):
\weq{8}{D_{z'} = z'^{2(r_1+r_2+\cdots+r_n)}D_z.}

If, then, $\delta$ is the degree of $D_z$, then
\[
2(r_1+r_2+\cdots+r_n) - \delta
\]
is the number of the vanishing roots of $D_{z'}$, and from \S\,190, 5. it follows:
\weq{9}{w_z = 2(r_1+r_2+\cdots+r_n).}

7. \gesperrt{The branching number $w_z$ is therefore always an even number.}

\begin{center}
\textbf{\S\,194. Differential quotients.}
\end{center}

Where we make no use of \gesperrt{continuity}, we cannot introduce the differential quotients in the ordinary way. We take two functions $\alpha, \beta$ from $\Omega$, between which the irreducible equation
\weq{1}{F(\alpha, \beta) = 0}
holds; we denote by $F'(\alpha), F'(\beta)$ the derived functions and \gesperrt{define} the differential quotient as a function of the field $\Omega$ by the formula:
\weq{2}{\left(\frac{d\alpha}{d\beta}\right) = -\frac{F'(\beta)}{F'(\alpha)}.}
%CONT>>
